\section{Prompt Optimization Overview}
\label{sec:overview}

\begin{figure*}[t]
  \centering
  \begin{tikzpicture}[
    font=\scriptsize,
    >={Latex[length=2.2mm]},
    flow/.style={->, semithick, draw=RRGray},
    box/.style={draw=RRGray, rounded corners=2pt, semithick, align=center,
      inner xsep=6pt, inner ysep=5pt},
    pool/.style={box, fill=RRGray!5, text width=2.35cm, minimum height=1.45cm},
    screen/.style={box, draw=RRBlue, fill=RRBlue!5,
      text width=3.15cm, minimum height=1.45cm},
    acquire/.style={box, draw=RROrange, fill=RROrange!7,
      text width=3.55cm, minimum height=1.55cm},
    data/.style={box, draw=RRBlue, fill=RRBlue!5,
      text width=3.15cm, minimum height=1.55cm},
    soft/.style={box, draw=RRGreen, fill=RRGreen!7,
      text width=4.2cm, minimum height=1.35cm},
    discrete/.style={box, draw=RRBlue, fill=RRBlue!4,
      text width=4.2cm, minimum height=1.35cm},
    deploy/.style={box, fill=RRGray!5, text width=5.8cm, minimum height=0.9cm}
  ]
    \node[pool] (pool) at (-0.1,1.7)
      {\textbf{1. Unlabeled pool}\\[2pt]
       $N$ candidate tuples $(x,q)$\\no gold answers};

    \node[screen] (screen) at (3.1,1.7)
      {\textbf{2. Paired screening}\\[2pt]
       run \textbf{Nano and Mini} on every tuple\\
       compare their predicted answers};

    \node[acquire] (acquire) at (7.25,1.7)
      {\textbf{3. Agreement-aware selection}\\[2pt]
       \textbf{disagree:} prioritize critical-case candidates\\
       \textbf{agree:} retain a coverage sample\\
       select only $B\ll N$ tuples; gold remains hidden};

    \node[data] (data) at (11.65,1.7)
      {\textbf{4. Annotate selected subset}\\[2pt]
       humans provide gold only for $B$ tuples\\
       derive \bc/\mo/\no/\bw, target, and $w=1/\bar c_y$};

    \node[soft] (soft) at (7.9,-0.65)
      {\textbf{5a. Primary: soft-prompt router}\\[2pt]
       frozen compact encoder\\train prompt vectors + linear head};

    \node[discrete] (discrete) at (13.2,-0.65)
      {\textbf{5b. Alternative: discrete router}\\[2pt]
       generative instruction router\\optimize policy text with GEPA};

    \node[deploy] (deploy) at (10.55,-2.55)
      {\textbf{6. Deploy one pre-router}\\
       $(x,q) \rightarrow \{\text{Nano},\text{Mini}\} \rightarrow$ operator output};

    \draw[flow] (pool) -- (screen);
    \draw[flow] (screen) -- node[above, text=RRGray, font=\tiny]
      {agree / disagree} (acquire);
    \draw[flow] (acquire) -- node[above, text=RRGray, font=\tiny]
      {$B$ selected IDs} (data);
    \coordinate (split) at (10.6,0.35);
    \draw[semithick, draw=RRGray] (data.south) -- (split);
    \fill[RRGray] (split) circle (1.1pt);
    \draw[flow] (split) -| (soft.north);
    \draw[flow] (split) -| (discrete.north);
    \draw[flow] (soft.south) |- (deploy.west);
    \draw[flow] (discrete.south) |- (deploy.east);

    \node[font=\tiny\bfseries, text=RRGray, anchor=south west]
      at (-1.55,2.55) {OFFLINE ACQUISITION};
    \node[font=\tiny\bfseries, text=RRGray, anchor=south west]
      at (5.65,0.18) {INDEPENDENT TRAINING PATHS};
  \end{tikzpicture}
  \caption{End-to-end workflow for an initially unlabeled candidate pool. Nano
  and Mini first screen every candidate without gold labels. Active selection
  prioritizes their disagreements, where recoverable one-model-only cases can
  occur, while retaining agreement examples for coverage. Humans annotate only
  the selected subset; the resulting gold-labeled paired set trains either the
  primary lightweight soft-prompt router or the independent discrete router.
  The two routers are alternatives, not stages of one model.}
  \Description{An unlabeled candidate pool is screened by Nano and Mini. Their
  predicted answers partition the pool into agreement and disagreement strata.
  Active selection prioritizes disagreements and samples some agreements, after
  which humans provide gold labels only for the selected subset. Gold labels and
  the cached paired predictions determine the four correctness outcomes and the
  cost-weighted route targets. This training set branches to two independent
  paths, a lightweight soft-prompt router and a discrete prompt router. Either
  trained router can choose Nano or Mini at deployment time.}
  \label{fig:overview}
\end{figure*}

Figure~\ref{fig:overview} gives the end-to-end workflow when the candidate pool
has no gold answers. The answer-model pair is fixed: $M_N$ denotes the cheap
Nano model and $M_M$ the stronger Mini model. We first execute both models on
the unlabeled pool and partition candidates using the observable event
$\hat z_N(x,q)=\hat z_M(x,q)$. Disagreements receive higher acquisition priority
because every Mini-only and Nano-only case must lie in this stratum; a sampled
agreement subset preserves coverage of ordinary cases and helps control
over-escalation. Selection never inspects an unselected gold answer. Humans then
annotate only the selected $B\ll N$ tuples, at which point the cached pair of
predictions and gold label reveal the \bc/\mo/\no/\bw outcome, route target, and
inverse-cost weight. The resulting agreement-stratified paired set is shared by
two independent training paths: the primary lightweight soft-prompt router and
the discrete prompt router. This workflow saves $N-B$ human annotations, though
it deliberately pays for paired model screening on the full candidate pool.

\paragraph{Soft-prompt router.}
The main \system router prepends $p$ learned embedding vectors to each tuple
and feeds the sequence to a frozen DistilBERT encoder
\cite{sanh2019distilbert}. Only the $p\times768$ prompt matrix and a two-class
linear head are trainable: 3,074, 4,610, 7,682, and 13,826 parameters for
$p\in\{2,4,8,16\}$. The frozen backbone still participates in inference, but
the task-specific state is tiny and routing is a single non-generative encoder
pass rather than autoregressive decoding by a billion-parameter LLM. This makes
the soft-prompt path the practical focus of \system.

Active acquisition determines which tuples train this router. The cost score
defines a hard route target, and its target route receives weight
$1/\bar c_y$. The soft prompt is trained only with the resulting weighted
cross-entropy (Eq.~\ref{eq:weighted-ce}). No outcome-specific constant is
tuned: disagreement enrichment, rather than a manual MO multiplier, ensures that
recoverable failures appear often enough to affect training. A
validation threshold exposes the cost--quality curve without retraining. The
design is inexpensive to retrain and supports multi-seed analysis, though its
learned prompt vectors are not human-readable
\cite{lester2021prompttuning}.
\paragraph{Independent discrete router.}
The discrete method is a separate router, not a component or second stage of
the soft-prompt method. It uses a Qwen3.5-2B generative execution router and a
Qwen3.5-4B reflection model during GEPA optimization. A seed instruction
defines the operator and route tokens; outcome-enriched demonstrations omit
gold labels and answer-model predictions. The adapter attaches route score and
textual feedback to each trace, and GEPA proposes revised policy text
\cite{agrawal2025gepa}. Reflection minibatches contain four examples and the
budget is 20 metric calls.

Hard routing asks for exactly \texttt{nano} or \texttt{mini}; invalid text
falls back to Mini. The scored interface emits raw risk
$\tilde{s}_P(x)\in[0,100]$, which the parser clips and normalizes to
$s_P(x)\in[0,1]$. Validation predicts Mini when $s_P(x)\ge\tau$; reported
thresholds 0.30 and 0.50 therefore use the normalized scale.

\paragraph{Role of the discrete comparison.}
The comparison tests whether the two innovations depend on one router class.
Discrete instructions offer interpretable policies but require substantially
larger generative models and can ignore a verbal rule at generation time. Soft
prompts are compact and directly supervised but can overfit a small validation
set. We report them as two independent methods under shared data and scoring,
never as an ensemble.


